\chapter{简介}

这是 \LUATEX\ 的参考手册。我们不敢说它完备齐全，并且假定读者已经了解
\TEX{}，正如 ``The \TEX\ Book、``\ETEX\ 手册、
``\PDFTEX\ 手册等文献中所描述的那样。更多参考资料发表于用户
组期刊以及与 \CONTEXT\ 相关的文档中。

从项目启动到稳定的 1.0 版本，大约经历了十年时间，但这是有充分理由的。
后续版本带来了新功能、更强的控制力以及对内部结构的清理。实验性
功能逐步演变为稳定功能，或被移除。早在相当早的阶段，\LUATEX\ 就已经可
以用于生产环境，作者们也在日常工作中使用它。后续版本有时需要对
\LUA\ 接口进行调整，但核心概念保持不变。当前版本在功能上可以认为是稳
定的，不会再有根本性的变化。当然，我们之后可能会决定推出具有不同特性
的 2.00 版本。

不要指望 \LUATEX\ 的行为与 \PDFTEX\ 完全一致！虽然那个 8 位引擎的核心
功能是我们的起点，但它已经与 \OMEGA{}（\ALEPH）的方向性支持相结合。而
且，\LUATEX\ 由于采用宽字符（32 位）、众多寄存器以及大内存支持，行为
可能有所不同。生成的 \PDF\ 代码与 \PDFTEX\ 不同，但用户通常不会察觉。
它原生支持 \UTF\ 输入、支持大型（超过 8 位）字体，并且数学排版机制针
对 \OPENTYPE\ 数学做了调优。同时也支持方向性排版。日志输出可能与其他
引擎不同，而且随着我们的推进，差异可能会更大。当你运行纯 \TEX\ 时，
\LUATEX\ 无疑比 \PDFTEX\ 慢，但当你运行例如 \CONTEXT\ \MKIV\ 时，在
很多情况下它反而更快，尤其是当你有稍微复杂一些的文档或输入时。无论如
何，全链路 32 位加上更多功能是有代价的，但在现代机器上这已经不是真正
的问题。

测试是在 \CONTEXT\ 下进行的，但 \LUATEX\ 也应该能与其他宏包良好配合。
为此我们提供了通用字体处理器，这些处理器与 \CONTEXT\ 中使用的大体相同。
讨论具体的实现细节超出了本手册的范围。即使我们努力让 \LUATEX\ 保持精
干高效，这里已经有足够多的内容值得讨论。

\LUATEX\ 由若干相互关联但（目前仍）可区分的部分组成。源代码的组织方式
经过了调整，以便将所有这些组件粘合在一起。我们持续清理 \TEX\ 引擎中积
累的代码所带来的副作用（尤其是不再需要的代码）。

\begin{itemize}
    \item 我们最初采用了 \PDFTEX\ 1.40.9 版本的大部分代码。代码库被转换为
        \CCODE\ 并拆分为多个模块。实验性功能被移除，实用宏也未被继承，
        因为它们的功能可以用 \LUA\ 编程实现。后端接口命令的数量已减少
        到少数几个。所谓的扩展功能与核心分离（我们尽量让核心保持与原始
        \TEX\ 核心接近）。一些机制如展开和凸排（protrusion）由于清理和
        优化，行为可能与原始版本不同。一些基于 whatsit 的功能（图像支持
        和可复用内容）现在已成为核心功能。我们不再与 \PDFTEX\ 的开发同
        步。
    
    \item 包含了 \ALEPH\ RC4（源自 \OMEGA{}）的方向模型。相关的原始命令是
        \LUATEX\ 核心的一部分，但在节点级别，方向性支持不再基于所谓的
        whatsit，而是基于具有相关属性的真实节点。方向的数量被限制在有
        用的集合内，后端也已具备方向感知能力。
    
    \item \ALEPH\ 的 I/O 转换过程、tcx 文件以及 \ENCTEX\ 均不可用。这些
        与编码相关的功能已被基于 \LUA\ 的方案（reader 回调）所取代。
        以类似的方式，所有文件 I/O\ 都可以被拦截。
    
    \item 我们目前使用 \LUA\ 5.3.*。有少数几个 \LUA\ 库我们认为是核心
        \LUA\ 机制的一部分，例如 \type {lpeg}。还有一些额外的 \LUA\
        库用于与 \TEX\ 的内部接口对接。我们也保留了 \LUA\ 5.2 的
        \type {bit32} 库。
    
    \item 有各种各样的 \TEX\ 扩展，但仅限于那些无法通过 \LUA\ 接口实现的
        扩展。数学机制通常有两条代码路径：一条是传统路径，另一条更适合
        宽 \OPENTYPE\ 字体。在这方面我们尽可能遵循 \MICROSOFT\ 规范。一
        些数学功能已经开放了一些接口，以便用户获得更多控制权。
    
    \item fontloader 使用了 \FONTFORGE\ 2008.11.17 的部分代码，结合了专为
        \TEX\ 引擎使用而编写的额外代码。我们尽量将特定字体支持最小化
        到 \TEX\ 所需的程度：字符引用和尺寸，并将其他一切委托给
        \LUA{}。这样我们可以在不触及核心的情况下让 \TEX\ 保持可扩展。
        为了最小化依赖，我们可能会在某个时候决定将其设为可选库。
    
    \item \METAPOST\ 库是 \LUATEX\ 的组成部分。这为 \TEX\ 提供了一定的
        图形能力，使用的是相对高速的图形子系统。同样，\LUA\ 被用作前端
        和后端之间的桥梁。\METAPOST\ 的进一步开发与 \LUATEX\ 密切相关。
    
    \item \TEX\ 自带的虚拟字体技术已被集成到字体机制中，允许在运行时创建
        虚拟字体。由于 \LUATEX\ 同时也可以充当 \LUA\ 解释器，这意味着
        无需额外程序即可构建完整的 \TEX\ 工作流。
    
    \item 从 1.09 版本开始，不再使用 poppler 库来包含（PDF），而是使用
        一个轻量级的专用库。这移除了一个依赖，但也使得 \LUATEX\ 的包
        含代码与 \PDFTEX\ 不同。事实上，由于 \LUA\ 图像接口的存在，它
        本来就已经大不相同了。
    
\end{itemize}

我们努力保持后续版本的兼容性，但中间版本可能包含实验性功能。一般规则
是：进入 \TEXLIVE\ 以及/或者在 \CONTEXT\ 会议前后发布的版本是稳定版
本。年度 \TEXLIVE\ 发布之间的任何版本都应视为 beta 版，在代码仓库中以
trunk 版本的形式存在。我们有一个用于开发的实验分支，但不对其任何实验
性功能提供支持。中间版本（来自 trunk）通常通过 \CONTEXT\ 发行渠道
（garden 以及所谓的 minimals）提供。

1.10 版本或多或少是开发的终点：这就是你所得到的。因为不仅 \CONTEXT\
（我们可以相当容易地适配）使用 \LUATEX{}，所以我们不能在没有不可预见后
果的情况下改变基础。到目前为止，已经证明 \LUA\ 可以用来扩展核心功能，
因此没有必要再添加更多功能，更不用说为（不那么）常见的问题提供硬编码解
决方案了。当然，会有 bug 修复，也许还有一些优化，甚至可能有一些新增内
容或非侵入式改进，但只能在稳定版本之外经过测试后才能加入。毕竟，二进制
文件已经够大了，能获得的收益也不多。

你可能会发现一些尚未文档化的 \LUA\ 辅助函数。这些被认为是实验性的，
不过当你在一个已经存在一段时间的 \CONTEXT\ 版本中遇到它们时，可以认
为它们会保留下来。当然，也可能只是我们还没来得及文档化它们。

这样的手册并不真正适合作为教程，关于教程我们推荐参阅 \CONTEXT\ 附带
的文档、文章等。它对所有读者来说也永远不够完备。我们努力跟上进度，但
读者需要意识到，这都是业余时间做的志愿者工作。而且可以肯定的是，抱怨手
册不好或文档糟糕并不会真正激励我们花更多时间在上面。话虽如此，我们希望
本文档是有用的。

\medskip

\needspace{5\baselineskip}

\begin{flushleft}
Hans Hagen
Harmut Henkel
Taco Hoekwater
Luigi Scarso
\end{flushleft}

\medskip

\begin{longtable}{l l}
版本     & \currentdate \\
\LUATEX  & 1.10 \\
\CONTEXT & MkIV \\
\end{longtable}
